\chapter{Introduction}\label{ch:introduction}



\section{Epilepsy and Seizure Prediction}
\textbf{Epilepsy} is a chronic neurological disorder characterized by recurrent, transient episodes of abnormal, synchronous neuronal activity known as seizures.
Globally, epilepsy affects approximately 65 million individuals and imposes a substantial social and economic burden, including increased rates of disability and mortality \parencite{leonardi_global_2002, feigin_global_2021}.

For the roughly one-third of patients with pharmacoresistant epilepsy, the \textbf{unpredictability of seizure occurrence} presents a significant challenge in daily life \parencite{schulze-bonhage_unpredictability_2008}.
They do not experience seizures more than \SI{99.95}{\percent} of the time, during which they would be able to live relatively normal lives; however, the constant uncertainty of seizure occurrences intrudes into every aspect of daily life \parencite{freestone_forward-looking_2017}.

Although some patients report clinical prodromi, such as behavioral or mood changes, which can alarm them to incoming seizures, the development of more reliable, automated \textbf{seizure prediction algorithms} holds the potential to improve patient quality of life by enabling personalized interventions, such as adaptive medication regimens, closed-loop neurostimulation, or timely precautionary actions by patients themselves, thereby increasing patient empowerment and reducing uncertainty \parencite{mormann_seizure_2007}.

Prediction requires data that contains seizure-related patterns and various \textbf{recording modalities} have been used for its acquisition.
These include non-invasive sensors for \gls{emg}, accelerometry, \gls{ekg}, \gls{ppg}, \gls{eda}, and skin temperature in various forms \parencite{brinkmann_seizure_2021}.
However, \gls{eeg} has widely been established to be the best modality, due to its exquisite temporal sensitivity which enables providing a near real-time display of ongoing cerebral activity \parencite{st_louis_electroencephalography_2016, freestone_forward-looking_2017}.
It can capture measurable shifts in brain dynamics hours or even days before a clinical event in many patients \parencite{cook_prediction_2013} and is discussed in further detail in the following section.



\section{Electroencephalography (EEG)}
\textbf{\Gls{eeg}} is an electrophysiological technique used to record the electrical activity arising from brains.
It primarily measures the extracellular ion flow generated by similarly oriented groups of cerebral cortical neurons---especially pyramidal neurons---near the scalp where the recording electrodes are placed in an area of roughly \SI{6}{\cm\squared} \parencite{st_louis_electroencephalography_2016}.
The measurement represents the \textbf{summation} of the \textbf{inhibitory} and \textbf{excitatory postsynaptic potentials}---rather than action potentials, which are too fast to measure---of neurons firing synchronously.
Because these neurons are oriented perpendicularly to the brain's surface, their combined activity creates a dipole that can be detected by recording electrodes.


\subsection{Clinical Utility and Applications}
The primary utility of \gls{eeg} lies in the evaluation of \textbf{dynamic cerebral functioning}. 
It is widely used for several clinical indications \parencite{st_louis_electroencephalography_2016}:
\begin{itemize}
    \item \textbf{Epilepsy:} \gls{eeg} is the gold standard for diagnosing and confirming seizures and epilepsy. 
    It allows clinicians to identify characteristic ictal (during-seizure) alterations and \glspl{ied}, such as spikes and sharp waves, which help classify seizure types and epilepsy syndromes.
    \item \textbf{General monitoring:} It is used to monitor the \textit{depth of anesthesia} during surgery and to detect sudden changes in neural functioning that may indicate \textit{ischemia or infarction}.
    \item \textbf{Research and specialized diagnosis:} \gls{eeg} waveforms can be averaged to study \textit{\glspl{ep}} and \textit{\glspl{erp}}, which represent neural activity related to specific stimuli, assisting in the analysis of cognitive and sensory functioning.
\end{itemize}

The following subsections discuss properties and modalities of \gls{eeg} that are important in the field of epilepsy.


\subsection{Frequency Bands}
\gls{eeg} activity is conventionally categorized into frequency\footnote{The exact frequency boundaries of each band vary in the literature and are based on \textcite{mormann_seizure_2007} here.} bandwidths, or bands (\cref{fig:eeg-bands}), which vary according to the patient's state (e.g., wakefulness, sleep) and developmental stage \parencite{st_louis_electroencephalography_2016, mormann_seizure_2007, nayak_normal_2026}:

{
\sisetup{range-phrase = {--}, range-units = single}
\begin{itemize}
    \item \textbf{Delta $\delta$} (\SIrange{0.5}{4}{\hertz}): Predominant during deep sleep and considered pathological if present in awake adults.
    \item \textbf{Theta $\theta$} (\SIrange{4}{8}{\hertz}): Commonly seen during drowsiness and in early sleep stages.
    \item \textbf{Alpha $\alpha$} (\SIrange{8}{12}{\hertz}): Present during relaxed, awake states, especially in the occipital head region. Strengthens when eyes are closed and attenuates during intense activity or concentration.
    \item \textbf{Beta $\beta$} (\SIrange{12}{35}{\hertz}): Associated with alert wakefulness and often enhanced by most sedative medications.
    \item \textbf{Gamma $\gamma$} (\SIrange{35}{100}{\hertz}): Associated with active cognition. Is not commonly seen in scalp \gls{eeg} but is readily detected in intracranial recordings.
\end{itemize}
}

\begin{figure}[ht]
    \centering
    \includegraphics[width=0.8\linewidth]{EEG_bands}
    \caption{The Raw EEG Signal and Corresponding Frequency Bands}
    \label{fig:eeg-bands}
\end{figure}


\subsection{EEG Recording Modalities}
The choice of electrode placement significantly impacts the recorded signal quality and spatial resolution \parencite{st_louis_electroencephalography_2016}:
\begin{itemize}
    \item \textbf{Surface (scalp) \gls{eeg}:} Non-invasive electrodes are placed on the skin using systems like the \textit{International 10--20 System} \parencite{klem_ten-twenty_1999}.
    However, the signal must traverse biological filters (brain, cerebrospinal fluid, meninges, skull, and skin), which reduce its amplitude and blur spatial information. 
    The skull acts as a \textit{spatial averager} and a dominant attenuator of high-frequency signals.
    \item \textbf{Video-\gls{eeg} Monitoring}: Considered the \enquote{gold standard} for characterizing paroxysmal spells and localizing seizures, this modality involves prolonged \gls{eeg} recording with simultaneous time-linked video. 
    Its primary value lies in the ability to correlate clinical behavioral changes (seizure semiology) with underlying electrographic activity. 
    It is essential for distinguishing true epileptic seizures from nonepileptic mimickers, such as psychogenic nonepileptic spells or syncope. 
    This monitoring is typically conducted in a specialized inpatient \gls{emu}, where antiepileptic drugs may be safely withdrawn to capture habitual seizures.
    \item \textbf{\Gls{ecog} / \gls{ieeg}:} Invasive monitoring involves placing electrodes directly on the cortex (subdural grids or strips) or within the brain (depth electrodes) \parencite{jayakar_diagnostic_2016}. 
    These provide much cleaner fidelity without muscle or movement artifacts, enabling precise localization of epileptic foci.
    \item \textbf{Subscalp / \gls{sqeeg}:} This modality involves electrodes implanted under the skin but outside the skull, such as the \textit{UNEEG SubQ system} used in this study.
    Such systems offer a robust and unobtrusive alternative to traditional scalp \gls{eeg} \parencite{duun-henriksen_subdural_2013}.
\end{itemize}

Research comparing subcutaneous systems to standard surface electrodes shows that they achieve similar performance, as measured by \textit{\gls{psd}} and \textit{coherence analysis}.
Subcutaneous recordings often demonstrate \textit{higher signal quality} during rest, primarily because the electrodes provide \textit{higher signal amplitudes} and maintain a \textit{rigid connection} to the source, reducing artifacts caused by electrode-skin instability or movements \parencite{duun-henriksen_eeg_2015}.
While the skull still acts as a spatial averager for these signals, the subcutaneous approach avoids degradation of electrode--skin contact over time (common with surface electrodes) and is portable and inconspicuous, making it suitable for \enquote{ultra-long-term} monitoring over months.



\section{Seizure-Related Intervals in EEG Recordings}\label{subsec:background-periictal-intervals-in-eeg-recordings}

In the analysis of \gls{eeg} data for seizure prediction, the continuous recording is typically categorized into four distinct functional states or intervals based on their temporal relationship to a seizure event: ictal, preictal, postictal, and interictal \parencite{moghim_predicting_2014}:

\begin{itemize}
    \item \textbf{Ictal Interval}: This state represents the period of the actual seizure activity. It is characterized by sudden, abnormal, and synchronous neuronal discharges that typically manifest as high-amplitude, highly variable signal energy across EEG channels. While the duration varies, these events often last approximately three minutes \parencite{moghim_predicting_2014}.
    \item \textbf{Preictal Interval}: This is the critical window immediately preceding a seizure where the brain undergoes a transition in dynamics leading to a seizure arising.
    \item \textbf{Postictal Interval}: This period follows the seizure offset during which the brain recovers from the ictus. EEG signals during this state may show abnormal excitement or depressed activity as physiological functions return to baseline.
    \item \textbf{Interictal Interval}: This state encompasses all non-seizure data that is temporally distant from both preictal and postictal periods. It serves as the baseline for a patient's \enquote{normal} brain activity between events.
\end{itemize}

During the distinct intervals, electrophysiological changes in the brain affect the voltages recorded in the \gls{eeg} electrodes (\cref{fig:background-periictal-intervals}), thus, understanding the specific signal properties of these intervals is fundamental to developing predictive algorithms.
To distinguish these states, various quantitative descriptors, or features, extracted from the \gls{eeg} signal, are analyzed. 
The following section will dive dive into the theory of which precursors have been investigated and how seizures might arise.

\begin{figure}[!htbp]
    \centering
    \includegraphics[width=\linewidth]{ictal_intervals}
    \caption[Annotated Seizure-Related Intervals of an \gls{ieeg} recording]{Annotated Seizure-Related Intervals of an \gls{ieeg} recording. Source: \textcite{moghim_predicting_2014}}
    \label{fig:background-periictal-intervals}
\end{figure}



\section{Precursors and Mechanisms of Ictogenesis}
Understanding the mechanisms by which a \enquote{normal} brain transitions into a seizure---a process known as \textit{ictogenesis}---is fundamental to predictive modeling \parencite{mormann_seizure_2007}.

It has been shown that various \textbf{clinical symptoms} often occur before seizures, which might be usable as predictive markers. 
They include increased \textit{oxygen availability} \parencite{adelson_noninvasive_1999}, \textit{cerebral blood flow} \parencite{weinand_cerebral_1997, baumgartner_preictal_1998}, \textit{blood-oxygen-level-dependent signals} \parencite{federico_functional_2005}, and changes in \textit{heart rate} \parencite{delamont_changes_1999, novak_time-frequency_1999, kerem_forecasting_2005}.

On a neural network level, one hypothesis for ictogenesis is that focal seizures arise from abnormally discharging \textit{\enquote{burster} neurons} that recruit and entrain neighbors into a \enquote{critical mass}.
The increasing intensity of critical interactions over time may be caused by increasing synchronization due to loss of inhibition, or by a lowering of the threshold for neuronal excitation which facilitates ictogenesis.
However, while these large-scale events should be easily accessible to conventional \gls{eeg} analysis, researchers have not detected it so far \parencite{mormann_seizure_2007, moghim_predicting_2014}.

Other groups have tried to explain ictogenesis using systems theory \parencite{maturana_critical_2020}.
With this paradigm, a key indicator of this transition is \enquote{\textit{critical slowing}}, which refers to the tendency of a system's regulatory mechanisms to slow down, thereby making the system take longer to return to equilibrium after perturbations.
This is viewed as a universal signature of a system losing its resilience as it approaches a tipping point or state transition.
In the context of epilepsy, this is often characterized by an increase in the autocorrelation of the \gls{eeg} or a lengthening of electrical evoked potentials \parencite{freestone_forward-looking_2017, maturana_critical_2020}.

This transition is frequently framed through the lens of bi/multistability, where a stable healthy attractor (e.g., wakefulness or sleep) coexists with a pathological seizure state.
Seizures may occur due to deterministic variations in underlying parameters or random \enquote{noisy} fluctuations in background activity that drive the brain's state across the energy barrier separating these attractors.
Because the brain possesses internal \textit{homeostatic regulatory mechanisms}, it is possible for it to enter a \textit{high susceptible} state without a clinical event actually occurring, as these mechanisms may correct the abnormal activity \parencite{freestone_forward-looking_2017}.
This conceptualization, among other lines of research, lead to the idea of a \textit{proictal state}---a period of heightened seizure likelihood that does not always terminate in an ictal event \parencite{muller_coherent_2022}.

While research into biomarkers like \textit{interictal spikes} and \textit{\glspl{hfo}} has revealed that while these indicators have clear implications for understanding ictogenesis, their dynamics are highly patient-specific \parencite{freestone_forward-looking_2017}.
Consequently, there is an increasing recognition that \enquote{one-size-fits-all} algorithms are insufficient, necessitating long-term monitoring to account for individual variability in seizure-generating mechanisms and the diverse populations of seizures that can coexist within a single patient \parencite{freestone_forward-looking_2017, yang_seizure_2024}.
The following section dives into further details about the history and methods of seizure prediction and forecasting.



\section{Seizure Prediction and its Evolution}
Using the aforementioned preictal interval, one paradigm for seizure prediction is to view it as a binary classification problem with the goal of distinguishing \textit{preictal} states from \textit{non-preictal}---especially interictal---states \parencite{yang_seizure_2024}. 
The following paragraphs will present the history of how this has been attempted.

The pursuit of identifying seizure precursors from the electroencephalogram began in the 1970s with \textit{linear approaches}, such as autoregressive modeling and power-spectrum analysis of spike-wave bursts \parencite{mormann_seizure_2007}.
Early research also focused on the predictive value of \textit{interictal spike rates}, though results were often contradictory; while some early studies suggested focal spiking decreased before seizures, subsequent investigations on more extended databases showed no systematic or reproducible changes \parencite{mormann_seizure_2007}.

In the 1990s, the field shifted toward \textit{nonlinear dynamics}, treating the brain as a chaotic system \parencite{mormann_seizure_2007, kuhlmann_epilepsyecosystemorg_2018}.
Researchers utilized measures like the largest Lyapunov exponent to identify a decrease in chaoticity in the minutes preceding a seizure, alongside the dynamical similarity index and correlation density to characterize a transitional preictal phase \parencite{mormann_seizure_2007}.
This era was marked by significant optimism, with many studies reporting high sensitivity on small, highly selected datasets \parencite{mormann_seizure_2007, muller_coherent_2022}.
However, a period of \textit{skepticism} followed in the early 2000s as more rigorous evaluations on unselected, continuous multi-day recordings failed to reproduce these optimistic findings, highlighting issues with overfitting, lack of specificity evaluation, and the absence of out-of-sample statistical validation \parencite{mormann_seizure_2007, freestone_forward-looking_2017}.

To overcome these reproducibility challenges and address the \enquote{ceiling effect} in predictive performance, the field turned to \textit{crowdsourcing} and \textit{open competitions} using standardized, high-quality datasets obtained using the NeuroVista device, the first implantable seizure prediction system \parencite{cook_prediction_2013}. 
Platforms like \textit{Kaggle} hosted the American Epilepsy Society Seizure Prediction Challenge, providing \gls{ieeg} recordings from both humans and canines \parencite{brinkmann_crowdsourcing_2016}.
Subsequent efforts by \textit{epilepsyecosystem.org} continued this trend, enabling researchers to compare algorithms on ultra-long-term human recordings \parencite{kuhlmann_epilepsyecosystemorg_2018}.
These competitions demonstrated that \textit{machine learning} could achieve significantly better-than-chance performance in both humans and canines \parencite{brinkmann_crowdsourcing_2016}.
Additionally, the first long-term in-human seizure advisory system (NeuroVista) provided evidence that, in some cases, seizures can be predicted hours in advance \parencite{cook_prediction_2013}.

Winning solutions of the competitions typically utilized a variety of \enquote{hand-crafted} features---including band power spectra, statistical moments (skewness, kurtosis), and cross-channel linear correlations in both time and frequency domains---classified by \textit{Support Vector Machines (SVM)}, \textit{Random Forests}, or \textit{k-nearest neighbors} \parencite{kuhlmann_epilepsyecosystemorg_2018, muller_coherent_2022}.

More recently, \textit{deep learning} has emerged as a dominant paradigm, with \textit{\gls{cnn}} and \textit{\gls{lstm}} networks used to extract features directly from raw data, bypassing the need for manual feature engineering \parencite{kuhlmann_epilepsyecosystemorg_2018, muller_coherent_2022}.
Looking forward, \textcite{freestone_forward-looking_2017} suggest that the field is shifting from a deterministic preictal process underlying the seizure prediction paradigm toward a probabilistic proictal state, thereby accepting the inherent uncertainty and moving toward seizure forecasting. 



\section{Epileptic Cycles}
The shift toward seizure forecasting was initiated following several seminal studies \parencite{baud_multi-day_2018, proix_forecasting_2021} that identified a relationship between circadian (daily), multidien (multi-day), and circannual (yearly) cycles, and the occurrence of epileptic seizures. 
Seizures were found to preferentially occur at specific phases of these cycles, thereby establishing prolonged periods of elevated seizure propensity---the hypothesized \textbf{proictal} state---rather than direct seizure precursors, i.e., the preictal state. 
The following section dives into further detail regarding the intricacies of the aforementioned epileptic cycles.


\subsection{History}
Observations of temporal organization in epilepsy date back to antiquity, with Babylonian writings noting that seizures often occurred \enquote{according to the calendar} \parencite{wilson_translation_1990}. 
In the 18th century, clinical accounts debated whether these patterns were driven by external factors, such as lunar cycles, or endogenous clock-like mechanisms \parencite{mead_treatise_1748, karoly_cycles_2021}. 
Systematic study began in the late 19th and early 20th centuries within \enquote{epilepsy colonies}---institutionalized settings where controlled environments allowed for meticulous long-term charting. 
There, researchers identified distinct circadian chronotypes, classifying patients by their tendency toward nocturnal, diurnal, or mixed seizure patterns \parencite{gowers_epilepsy_1881, langdon-down_time_1929}. 
Furthermore, \textcite{griffiths_rhythm_1938} provided seminal evidence for multidien cycles, observing seizure periodicities ranging from one week to several months in both children and adults \parencite{karoly_cycles_2021, baud_multi-day_2018}.


\subsection{Modern Views on Cycles}
Modern neuroscience has validated these historical observations using ultra-long-term \gls{eeg} recorded via implanted devices. 
A central biomarker in this research is \gls{iea} (\cref{fig:iea}), which refers to transient electrographic discharges---such as spike-waves, polyspikes, and fast oscillations---that occur between seizures without producing overt clinical symptoms \parencite{baud_multi-day_2018, karoly_cycles_2021}. \Gls{iea} is typically quantified as counts of \glspl{ied} (\cref{fig:ied}) per unit of time and is regarded as a measure of \enquote{brain irritability} \parencite{karoly_cycles_2021}.

\begin{figure}[ht]
    \centering
    \begin{subfigure}[t]{0.45\textwidth}
        \centering
        \includegraphics[width=\textwidth]{ied}
        \caption{A Single \gls{ied}}
        \label{fig:ied}
    \end{subfigure}
    \hfill
    \begin{subfigure}[t]{0.54\textwidth}
        \centering
        \includegraphics[width=\textwidth]{iea}
        \caption{\Gls{iea}}
        \label{fig:iea}
    \end{subfigure}

    \caption{Examples of \Gls{iea} and \Glspl{ied} in the \Gls{eeg}. Source: \textcite{baud_multi-day_2018}}
    \label{fig:ied-and-iea-examples}
\end{figure}

Recent studies have demonstrated that \gls{iea} is not static but fluctuates according to robust circadian and multidien rhythms that are stable within individuals for up to a decade \parencite{baud_multi-day_2018}. While circadian \gls{iea} appears to be largely modulated by the sleep-wake cycle and associated hormones, such as melatonin, multidien rhythms appear to be endogenous and free-running, showing no clear alignment with universal environmental cues \parencite{karoly_cycles_2021}. 
Crucially, seizures are not randomly distributed across these cycles; they are significantly phase-locked, often occurring preferentially during the rising phase (up-slope) of multidien \gls{iea} rhythms \parencite{baud_multi-day_2018, karoly_cycles_2021}.


\subsection{Seizure Forecasting and the Proictal Hypothesis}
The discovery of these multi-scale rhythms has enabled the paradigm shift toward probabilistic seizure forecasting, an approach comparable to weather forecasting \parencite{proix_forecasting_2021, karoly_cycles_2021}. 
This framework overtly distinguishes between the preictal and proictal state. 
In a proictal state, the probability of a seizure is elevated due to the alignment of high-risk phases across multiple cycles, yet the actual occurrence remains stochastic; a patient may experience multiple seizures or none at all during a high-risk period \parencite{karoly_cycles_2021}.
As mentioned, it may be possible for the brain to enter a highly susceptible (proictal) state without epileptic events occurring, because its homeostatic, regulatory mechanisms may correct abnormal activity \parencite{buzsaki_rhythms_2006, freestone_forward-looking_2017}.
By leveraging the phase information of \gls{iea} cycles, researchers have attempted to forecast periods of high risk days in advance, providing a \enquote{seizure gauge} that replaces constant uncertainty with quantified probability \parencite{proix_forecasting_2021}.

However, in the context of seizure forecasting and prediction, two major challenges persist: (1) understanding if or why some seizures are not preceded by detectable preictal changes and thus remain unpredictable \textbf{(sensitivity)}, and (2) elucidating the origins of the number of false positive predictions \textbf{(specificity)}.
One hypothesis is that not all seizures are preceded by reliable preictal signatures, reflecting truly spontaneous ictogenesis.
Alternatively, limitations in model architecture or feature selection may prevent the detection of existing preictal patterns.
A further possibility, which is investigated in this work, is that false positive predictions may correspond to periods when the brain is in a proictal state, but the critical perturbation did not occur or a threshold was not surpassed, and thus the seizure never arose, thus making the prediction not truly \enquote{false} in a stochastic sense \parencite{muller_coherent_2022}.


\subsection{Coherent False Predictions}
A significant challenge in the field has been the \enquote{ceiling effect} of prediction performance, characterized by a high number of false positive alarms.
\textcite{muller_coherent_2022} investigated whether these errors resulted from algorithmic limitations or the nature of the data itself. 
They compared two substantially different classifiers: a \gls{cnn} processing raw \gls{ieeg} signals and \glspl{mlp} operating on hand-crafted features \parencite{muller_coherent_2022}.

They observed a high correlation of false predictions between these two fundamentally different methods. 
Since both the raw signal and feature-based classifier made false predictions in a coherent manner, this provided evidence that the errors were not caused by failing features but were rooted in the data. 
They interpreted these co-occurring false positives as evidence of the \enquote{proictal} state as an epoch of high seizure likelihood and ictogenic potential that did not culminate in an overt ictal event. 
Under this view, current binary classification models are limited because they assume any period not followed by a seizure is, by definition, \enquote{interictal,} potentially ignoring genuine states of increased brain irritability.



\section{Research Goals and Approach}
This work systematically investigated the hypothesis posed by \textcite{muller_coherent_2022} that \glspl{fp} are indeed manifestations of an underlying proictal state. 
To accomplish this, two predictive models were trained that are highly similar to the ones in \textcite{muller_coherent_2022}: a \gls{cnn} that operates on raw, subcutaneous \gls{eeg} signals and an ensemble of \glspl{mlp} that operate on hand-crafted features.
If \gls{fp} predictions appear preferentially in the same phases as true seizures (e.g., the up-slope of the multidien cycle), it would provide evidence that these predictions are not mere algorithmic \enquote{noise} but are correctly identifying proictal states that simply were not followed by a seizure.
Additionally, temporal co-occurrence of false positives across these two models was interpreted as evidence of a proictal state.

The \textit{primary goal} of this work was thus to provide a joint analysis of these two pillars of evidence to make progress in the field towards a more comprehensive model of seizure risk.
It was hypothesized that a \textit{co-occurrence of \glspl{fp}} would coincide with a \textit{similar phase preference of \glspl{fp} and seizures}.

The \textit{secondary goal} was to evaluate the model performance.
It was investigated how it compares to the existing literature, which is of particular interest due to the ultra-long-term recordings used here.
As in \textcite{muller_coherent_2022}, it was also investigated how the two model families compare to each other.
